\documentclass[11pt,twocolumn,oneside,openany,headings=optiontotoc,11pt,numbers=noenddot]{article}

\usepackage[a4paper]{geometry}
\usepackage[utf8]{inputenc}
\usepackage[T1]{fontenc}
\usepackage{lmodern}
\usepackage[ngerman]{babel}
\usepackage{ngerman}

\usepackage[onehalfspacing]{setspace}

\usepackage{fancyhdr}
\usepackage{fancybox}

\usepackage{rotating}
\usepackage{varwidth}

%Struktogramme
\usepackage[german,curves]{struktex}

\usepackage{pdflscape}
\usepackage{changepage}
\usepackage{graphicx}
\usepackage[bottom]{footmisc}
\usepackage{transparent}
\usepackage{graphbox}
\graphicspath{
	{Pics/PDFs/}
	{Pics/JPGs/}
	{Pics/PNGs/}
}
\usepackage{caption}
\usepackage{wrapfig}
\usepackage{marginnote}
\usepackage{tabularx}
\usepackage{dashrule}
\usepackage{soulutf8}
\usepackage{hhline}
%arydshln suppresses vertical lines in table
%\usepackage{arydshln}
\usepackage{multirow}
\usepackage{enumerate}
\usepackage[hidelinks]{hyperref}
\usepackage{listings}

\usepackage[table]{xcolor}
\usepackage{array}
\usepackage{enumitem,amssymb,amsmath}
\usepackage{interval}
\usepackage{cancel}
\usepackage{stmaryrd}
\usepackage{wasysym}
\usepackage{polynom}
\usepackage{diagbox}
\usepackage{dashrule}
\usepackage{framed}
\usepackage{mdframed}
\usepackage{karnaugh-map}
\usepackage{pdfpages}

\usepackage{blindtext}

\usepackage{eso-pic}

\usepackage{amssymb}
\usepackage{eurosym}

\usepackage[pages=some]{background}
\pagestyle{headings}
\renewcommand{\headrulewidth}{0.2pt}
\renewcommand{\footrulewidth}{0.2pt}
\newcommand*{\underdownarrow}[2]{\ensuremath{\underset{\overset{\Big\downarrow}{#2}}{#1}}}
\setlength{\fboxsep}{5pt}
\newcommand{\explainBelow}[3]{\underbrace{#1}_{\parbox{\widthof{#3}}{\footnotesize\raggedright #2}}}
\newcommand{\explainAbove}[3]{\overbrace{#1}^{\parbox{\widthof{#3}}{\footnotesize\raggedright #2}}}
\newcommand\footnoteref[1]{\protected@xdef\@thefnmark{\ref{#1}}\@footnotemark}


% Codestyle defined
\definecolor{codegreen}{rgb}{0,0.6,0}
\definecolor{codegray}{rgb}{0.5,0.5,0.5}
\definecolor{codepurple}{rgb}{0.58,0,0.82}
\definecolor{backcolour}{rgb}{0.95,0.95,0.92}
\definecolor{deepgreen}{rgb}{0,0.5,0}
\definecolor{darkblue}{rgb}{0,0,0.65}
\definecolor{mauve}{rgb}{0.40, 0.19,0.28}
\colorlet{exceptioncolour}{yellow!50!red}
\colorlet{commandcolour}{blue!60!black}
\colorlet{numpycolour}{blue!60!green}
\colorlet{specmethodcolour}{violet}

%Neue Spaltendefinition
\newcolumntype{L}[1]{>{\raggedright\let\newline\\\arraybackslash\hspace{0pt}}m{#1}}
\newcolumntype{M}{>{\centering\arraybackslash}X}
\newcommand{\cmnt}[1]{\ignorespaces}
%Textausrichtung ändern
\newcommand\tabrotate[1]{\rotatebox{90}{\raggedright#1\hspace{\tabcolsep}}}

%Intervall-Konfig
\intervalconfig {
	soft open fences
}

%Bash
\lstdefinestyle{BashInputStyle}{
	language=bash,
	basicstyle=\small\sffamily,
	backgroundcolor=\color{backcolour},
	columns=fullflexible,
	backgroundcolor=\color{backcolour},
	breaklines=true,
}
%Java
\lstdefinestyle{JavaInputStyle}{
	language=Java,
	backgroundcolor=\color{backcolour},
	aboveskip=1mm,
	belowskip=1mm,
	showstringspaces=false,
	columns=flexible,
	basicstyle={\footnotesize\ttfamily},
	numberstyle={\tiny},
	numbers=none,
	keywordstyle=\color{purple},,
	commentstyle=\color{deepgreen},
	stringstyle=\color{blue},
	emph={out},
	emphstyle=\color{darkblue},
	emph={[2]rand},
	emphstyle=[2]\color{specmethodcolour},
	breaklines=true,
	breakatwhitespace=true,
	tabsize=2,
}
%Python
\lstdefinestyle{PythonInputStyle}{
	language=Python,
	alsoletter={1234567890},
	aboveskip=1ex,
	basicstyle=\footnotesize,
	breaklines=true,
	breakatwhitespace= true,
	backgroundcolor=\color{backcolour},
	commentstyle=\color{red},
	otherkeywords={\ , \}, \{, \&,\|},
	emph={and,break,class,continue,def,yield,del,elif,else,%
		except,exec,finally,for,from,global,if,import,in,%
		lambda,not,or,pass,print,raise,return,try,while,assert},
	emphstyle=\color{exceptioncolour},
	emph={[2]True,False,None,min},
	emphstyle=[2]\color{specmethodcolour},
	emph={[3]object,type,isinstance,copy,deepcopy,zip,enumerate,reversed,list,len,dict,tuple,xrange,append,execfile,real,imag,reduce,str,repr},
	emphstyle=[3]\color{commandcolour},
	emph={[4]ode, fsolve, sqrt, exp, sin, cos, arccos, pi,  array, norm, solve, dot, arange, , isscalar, max, sum, flatten, shape, reshape, find, any, all, abs, plot, linspace, legend, quad, polyval,polyfit, hstack, concatenate,vstack,column_stack,empty,zeros,ones,rand,vander,grid,pcolor,eig,eigs,eigvals,svd,qr,tan,det,logspace,roll,mean,cumsum,cumprod,diff,vectorize,lstsq,cla,eye,xlabel,ylabel,squeeze},
	emphstyle=[4]\color{numpycolour},
	emph={[5]__init__,__add__,__mul__,__div__,__sub__,__call__,__getitem__,__setitem__,__eq__,__ne__,__nonzero__,__rmul__,__radd__,__repr__,__str__,__get__,__truediv__,__pow__,__name__,__future__,__all__},
	emphstyle=[5]\color{specmethodcolour},
	emph={[6]assert,range,yield},
	emphstyle=[6]\color{specmethodcolour}\bfseries,
	emph={[7]Exception,NameError,IndexError,SyntaxError,TypeError,ValueError,OverflowError,ZeroDivisionError,KeyboardInterrupt},
	emphstyle=[7]\color{specmethodcolour}\bfseries,
	emph={[8]taster,send,sendMail,capture,check,noMsg,go,move,switch,humTem,ventilate,buzz},
	emphstyle=[8]\color{blue},
	keywordstyle=\color{blue}\bfseries,
	rulecolor=\color{black!40},
	showstringspaces=false,
	stringstyle=\color{deepgreen}
}

\lstset{literate=%
	{Ö}{{\"O}}1
	{Ä}{{\"A}}1
	{Ü}{{\"U}}1
	{ß}{{\ss}}1
	{ü}{{\"u}}1
	{ä}{{\"a}}1
	{ö}{{\"o}}1
}

% Neue Klassenarbeits-Umgebung
\newenvironment{worksheet}[3]
% Begin-Bereich
{
	\newpage
	\sffamily
	\setcounter{page}{1}
	\ClearShipoutPicture
	\AddToShipoutPicture{
		\put(55,761){{
				\mbox{\parbox{385\unitlength}{\tiny \color{codegray}BBS I Mainz, #1 \newline #2
						\newline #3
					}
				}
			}
		}
		\put(455,761){{
				\mbox{\hspace{0.3cm}\includegraphics[width=0.2\textwidth]{../../logo.pdf}}
			}
		}
	}
}
% End-Bereich
{
	\clearpage
	\ClearShipoutPicture
}

\setlength{\columnsep}{3em}
\setlength{\columnseprule}{0.5pt}

\geometry{left=2.00cm,right=2.00cm,top=3.00cm,bottom=1.00cm,includeheadfoot}
\pagenumbering{gobble}
\pagestyle{empty}

\begin{document}
	\begin{worksheet}{HBF IT 17A}{Grundstufe}{Verhalten für große und kleine x-Werte | Differenzialrechnung}
		\section{Verhalten für kleine bzw. große x-Werte}
		Haben wir eine Funktion \[f(x) = a_nx^n + a_{n-1}x^{n-1} + \ldots + a_1x + a_0\] gegeben und möchten für diese das Verhalten für kleine bzw. große x-Werte beschreiben, so betrachten wir den charakteristischen Term \(\mathbf{a_nx^n}\).\\
		Für das Verhalten gilt folgendes\\
		\par\noindent
		\begin{tabular}{l|L{0.17\textwidth}|L{0.17\textwidth}}
			\hline
			\diagbox{\(a_{n}\)}{n} & gerade & ungerade\\
			\hline
			positiv & \shortstack{\(f(x) \xrightarrow{x \rightarrow -\infty} \infty\)\\\(f(x) \xrightarrow{x \rightarrow \infty} \infty\)} & \shortstack{\(f(x) \xrightarrow{x \rightarrow -\infty}-\infty\)\\\(f(x) \xrightarrow{x \rightarrow \infty}\infty\)}\\
			\hline
			negativ & \shortstack{\(f(x) \xrightarrow{x \rightarrow -\infty}-\infty\)\\\(f(x) \xrightarrow{x \rightarrow \infty}-\infty\)} & \shortstack{\(f(x) \xrightarrow{x \rightarrow -\infty}\infty\)\\\(f(x) \xrightarrow{x \rightarrow \infty}-\infty\)}\\
			\hline
		\end{tabular}
		\section{Rechnerisch die Nullstellen (NST) bestimmen}
		Abhängig von der gegebenen Funktion haben wir verschiedene Möglichkeiten die NST dieser zu bestimmen.
		\begin{itemize}
			\item \(a_1x + a_0 \Rightarrow\) Null setzen und nach \(x\) umformen (\(a_1x + a_0 = 0 \Rightarrow x = \frac{a_0}{a_1}\))
			\item  \(a_2x^2 + a_1x + a_0 \Rightarrow\) Anwenden der pq-Formel\\
			\textbf{Beachte:} \(a_2\) muss den Wert 1 haben (also muss gegebenenfalls \(:a_2\) gerechnet werden).
			\item \(a_2x^2 + a_0 \Rightarrow\) Null setzen und nach \(x\) umformen (\(a_2x^2 + a_0 = 0 \Rightarrow x = \sqrt{\frac{a_0}{a_2}}\))
			\item Funktionen mit Grad 3 und Höher \(\Rightarrow\) NST ausprobieren (meist \(-2,-1,0,1,2\)); Dann mit Polynomdivision den Grad reduzieren\\
			\((a_nx^n + \ldots +a_1x +a_0):(x-NST) = p(x)\)\\
			NST des Ergebnisses \(p(x)\) bestimmen
			\item Weitere Möglichkeit: Durch Substitution\\
			\(2x^4 -10x^2+8 \xrightarrow{z:=x^2} 2z^2-10z+8\)\\
			Hat man hiervon die NST bestimmt, muss Rücksubstituiert werden (also \(x^2 = NST \Rightarrow x = \pm \sqrt{NST}\))
		\end{itemize}
		\section{Zuordnen bzw. Skizzieren des Graphen der Ableitungsfunktion}
		Haben wir den Funktionsgraphen der Funktion \(f(x)\) gegeben und möchten den Graph der Ableitungsfunktion zeichnen, gehen wir wie folgt vor:\\
		\begin{itemize}
			\item Markiere die x-Werte der Hoch- bzw. Tiefpunkte auf der x-Achse
			\item Betrachte die Steigung des Funktionsgraphen links / zwischen / rechts von den markierten x-Werten
			\begin{itemize}
				\item Steigung positiv \(\Rightarrow\) markiere oberhalb der x-Achse
				\item Steigung negativ \(\Rightarrow\) markiere unterhalb der x-Achse
			\end{itemize}
			\item Skizziere den Graph der Ableitung entsprechend der Markierungen (oberhalb / unterhalb). Der Graph muss dabei durch die markierten x-Werte gehen.
		\end{itemize}
		Um den Ableitungsgraphen zuzuordnen, können wir ähnlich vorgehen. \(\Rightarrow\) Suche den Graphen der Ableitungsfunktion, der an den Hoch- bzw. Tiefpunkten des Funktionsgraphen die x-Achse schneidet (also eine NST besitzt).
		\section{Extremwerte rechnerisch bestimmen}
		Extremstellen einer Funktion sind eben solche Stellen, an denen die Steigung \(m\) des Funktionsgraphen von \(f(x)\) 0 ist. Dies ist genau dann der Fall, wenn die Ableitungsfunktion \(f'(x)\) den Wert 0 annimmt. Also wenn gilt: \(\mathbf{f'(x) = 0}\)\\
		Wir gehen also wie folgt vor, um die Extrempunkte zu bestimmen:
		\begin{itemize}
			\item Bestimme zunächst die Ableitungsfunktion \(f'(x)\)
			\item Berechne die Nullstellen (NST) der Ableitungsfunktion \(f'(x)\)
			\item Ermittle die Koordinaten der berechneten Extremstellen. Dafür setze die x-Koordinate der NST in die Ausgangsfunktion ein (\(f(NST)\)).
		\end{itemize}	
	\end{worksheet}
\end{document}